\section{Conclusions}
\label{sec:conclusion}

Eight search strategies were compared under an equal budget of 287 evaluations beyond a shared pool of 9,800 designs on the conceptual design of a 1.5\,m-span, 8\,kg electric tailless blended-wing-body UAV, and the evaluator and its stability constraint were audited against independent models. The conclusions are the following.

\begin{enumerate}
    \item \textbf{Search strategies.} TuRBO gave the longest range, $\NvMeanTurbo \pm \NvSdTurbo$\,km (\NvGainPctTurbo\% above the pool best). The neural-ensemble active-learning loop as originally implemented reached $\NvMeanAL \pm \NvSdAL$\,km and was not better than an adaptive local search without a surrogate ($\NvMeanAdapt \pm \NvSdAdapt$\,km, Welch $p \NvPALvAdapt$). With a repair that enforces the washout ordering before evaluation it reached $\NvMeanALR \pm \NvSdALR$\,km, \NvDALRvAdapt\,km above the adaptive local search (Holm $p \NvHALRvAdapt$) and \NvDTurbovALR\,km below TuRBO (Welch $p \NvPTurbovALR$, Holm $p \NvHTurbovALR$, interval $\NvCITurbovALR$\,km). Adaptive local search alone recovers \NvShareAdapt\% of the gain of TuRBO over the pool. The ordering of TuRBO, AL-R and the adaptive local search is the same from pools of 288, 1{,}000 and 3{,}000 designs; a pool of 1{,}000 designs costs TuRBO \NpTurboLossB\,km when it is a subsample of the full pool and \NpIndTurboVsFull\,km when it is generated independently, and with 288 pool designs every TuRBO run still exceeds the best of the full pool. At twice the budget the ordering is unchanged, but the lead of TuRBO over the adaptive local search falls from \NlTvAAD\ to \NlTvACD\,km.
    \item \textbf{Feasibility handling.} The classifier weight is necessary for the active-learning loop: without it no feasible design was found in 20 runs, with or without the repair. Its form matters little at the full budget, and adding it to TuRBO raised the feasible fraction but not the mean range. Constraints that can be satisfied by construction should be repaired in every strategy before a benchmark is interpreted.
    \item \textbf{Stability constraint.} Euler CFD on \NcalN\ designs shows that the lattice static margin is low by a nearly constant \NcalA\ ($R^2 = \NcalRsq$). The $0.25$ cap therefore corresponds to a lattice value of \NcalCapEquiv. On \NoosN\ further run champions that were not used in the fit, the calibration predicts the Euler margin with a root-mean-square error of \NoosRmse. Applying the calibrated constraint leaves the ranking of the strategies unchanged, but the best design of a minority of runs of most strategies, with and without a surrogate, lies beyond the calibrated limit.
    \item \textbf{Bounds and loft.} Among single-bound relaxations only the battery-mass limit has a measurable price (\NbpBmassGain\,km for $0.2$\,kg); the best designs are limited by the energy they may carry and by the minimum cruise speed. The chord clamp of the planform loft is a validity guard with no detectable effect on range, and the loft does not guarantee a monotone center-body chord.
    \item \textbf{Evaluator.} Panel refinement changes $L/D$ of the good designs by $0.6$ to $1.5\%$, and the ranking agrees with AeroBuildup ($\rho = 0.98$), but absolute drag differs between the models by a median factor of $1.46$ and no viscous three-dimensional solution or experiment is available. Ranges are therefore comparative.
\end{enumerate}

The audit of the code against the text changed the results: two baselines were corrected, an asymmetry between strategies was removed, and one geometric claim was withdrawn. The first implementations are retained beside the corrected ones. For studies of this kind the recommendation is to include a surrogate-free control with the same proposal structure as the surrogate loop, to report all seeds including failed runs, and to calibrate any constraint that the optimum presses against with a higher-fidelity model before the search.
